\section{Hai nguyên nhân cốt lõi}
\label{sec:nguyennhan}

Phương trình~\eqref{eq:central} biểu diễn gradient của tầng $l$ như một tích gồm
$L-l$ thừa số, mỗi thừa số có dạng
\[
  \mathbf{M}^{[k]}
  = \operatorname{diag}\!\left(\sigma'\!\left(\mathbf{z}^{[k-1]}\right)\right)
    \left(\mathbf{W}^{[k]}\right)^{\!\top}.
\]
Mỗi thừa số là tích của \emph{hai} đối tượng, và mỗi đối tượng đóng góp một
nguyên nhân độc lập: ma trận đường chéo mang thông tin về hàm kích hoạt, ma trận
chuyển vị mang thông tin về trọng số. Ta phân tích lần lượt rồi hợp nhất.

\subsection{Nguyên nhân 1 --- độ bão hoà của hàm kích hoạt}

\begin{infobox}{Chặn trên của đạo hàm sigmoid}
Với $\sigma(z) = \left(1+e^{-z}\right)^{-1}$, ta có
\[
  \sigma'(z) = \sigma(z)\bigl(1-\sigma(z)\bigr)
  \qquad\text{và}\qquad
  \max_{z\in\mathbb{R}}\sigma'(z) = \frac{1}{4},
\]
đạt được \emph{duy nhất} tại $z = 0$.
\end{infobox}

\noindent\textit{Chứng minh.}\;
Đạo hàm trực tiếp cho
$\sigma'(z) = e^{-z}\left(1+e^{-z}\right)^{-2}$. Viết lại theo $\sigma$:
\[
  \sigma'(z)
  = \frac{1}{1+e^{-z}}\cdot\frac{e^{-z}}{1+e^{-z}}
  = \sigma(z)\bigl(1-\sigma(z)\bigr).
\]
Đặt $u = \sigma(z)\in(0,1)$. Hàm $u\mapsto u(1-u)$ là parabol lõm xuống, đạt cực
đại tại $u = 1/2$ với giá trị $1/4$. Vì $\sigma$ là song ánh tăng nghiêm ngặt từ
$\mathbb{R}$ lên $(0,1)$, giá trị $u = 1/2$ ứng với đúng một điểm $z = 0$. Do đó
$\sigma'(z)\le 1/4$ với mọi $z$, dấu bằng chỉ tại gốc.\hfill$\square$

\medskip
Hệ quả trực tiếp: ma trận đường chéo trong mỗi thừa số có chuẩn phổ
$\left\|\operatorname{diag}(\sigma'(\mathbf{z}))\right\|_2
= \max_i \sigma'(z_i) \le 1/4$. Ngay cả trong trường hợp \emph{thuận lợi nhất}
--- mọi tiền kích hoạt đều đúng bằng $0$ --- mỗi tầng vẫn nhân gradient với một
hệ số không quá $1/4$. Nếu $\left\|\mathbf{W}^{[k]}\right\|_2 \approx 1$, biên độ
gradient co lại theo $4^{-L}$.

\medskip
\noindent Bảng suy giảm dưới đây cho thấy tốc độ của hiện tượng:

\begin{table}[H]
\centering
\begin{tabular}{@{}crl@{}}
\toprule
\rowcolor{rowLavender}
\textbf{Số tầng $L$} & \textbf{$0.25^{L}$} & \textbf{Diễn giải} \\
\midrule
$1$  & $2.500000\times10^{-1}$  & suy giảm nhẹ, chưa đáng ngại \\
$5$  & $9.765625\times10^{-4}$  & tầng đầu học chậm hơn tầng cuối $\sim10^3$ lần \\
$10$ & $9.536743\times10^{-7}$  & chỉ còn $8$ lần epsilon máy \texttt{float32} \\
$20$ & $9.094947\times10^{-13}$ & dưới ngưỡng phân giải của \texttt{float32} \\
$50$ & $7.888609\times10^{-31}$ & bằng không về mặt số học \\
\bottomrule
\end{tabular}

\end{table}

Con số ở dòng $L=10$ đáng được nhấn mạnh. Epsilon máy của \texttt{float32} là
$\varepsilon_{32} \approx 1.192093\times10^{-7}$. Giá trị $0.25^{10} \approx
9.536743\times10^{-7}$ chỉ lớn hơn $\varepsilon_{32}$ đúng $8$ lần. Nghĩa là với
một mạng sigmoid $10$ tầng huấn luyện ở độ chính xác đơn, gradient tầng đầu đã
tiệm cận giới hạn biểu diễn của kiểu số: phép cộng $w \leftarrow w - \eta g$ với
$g$ nhỏ như vậy bị làm tròn về chính $w$, và tầng đó \emph{đóng băng hoàn toàn}.
Ở $L=20$, giá trị $9.094947\times10^{-13}$ đã nằm sâu dưới $\varepsilon_{32}$.

\subsubsection{Hàm tanh làm chậm vấn đề nhưng không loại bỏ nó}

Với $\tanh$, ta có $\tanh'(z) = 1 - \tanh^2(z)$, và
\[
  \max_{z\in\mathbb{R}}\tanh'(z) = 1,
\]
đạt được duy nhất tại $z = 0$. Chặn trên bằng $1$ thay vì $1/4$ là một cải thiện
thực chất: tại chế độ tuyến tính quanh gốc, $\tanh$ bảo toàn biên độ gradient
trong khi sigmoid làm co lại bốn lần mỗi tầng.

Tuy nhiên cải thiện này chỉ tồn tại \emph{tại một điểm duy nhất}. Với $|z|\ge 2$
đã có $\tanh'(z) \le 0.072$, nhỏ hơn cả chặn $1/4$ của sigmoid. Trong một mạng
thật, các tiền kích hoạt phân tán trên một khoảng chứ không tập trung tại $0$, nên
giá trị $\tanh'$ \emph{trung bình} luôn nhỏ hơn $1$ ngặt. Tích của $L$ số nhỏ hơn
$1$ ngặt vẫn hội tụ về $0$ theo cấp số nhân --- chỉ với cơ số lớn hơn.

\begin{notebox}
$\tanh$ dời ngưỡng chiều sâu mà vấn đề trở nên nghiêm trọng, chứ không thay đổi
bản chất cấp số nhân của nó. Đây là lý do việc chuyển từ sigmoid sang $\tanh$
giúp huấn luyện mạng $10$ tầng nhưng vẫn thất bại ở mạng $50$ tầng.
\end{notebox}

\begin{figure}[H]
\centering
\begin{tikzpicture}
\begin{axis}[
    width=0.92\linewidth, height=6.4cm,
    xlabel={$z$}, ylabel={giá trị},
    xmin=-6, xmax=6, ymin=-1.15, ymax=1.35,
    domain=-6:6, samples=201,
    axis lines=left,
    legend style={font=\scriptsize, at={(0.5,1.30)}, anchor=north,
                  legend columns=2, draw=none, fill=none},
    tick label style={font=\scriptsize},
    label style={font=\scriptsize},
    grid=major, grid style={gray!20},
]
% ---- saturation regions ----
\fill[gray!18] (axis cs:-6,-1.15) rectangle (axis cs:-3,1.35);
\fill[gray!18] (axis cs:3,-1.15)  rectangle (axis cs:6,1.35);
\node[font=\scriptsize, text=gray!60!black] at (axis cs:-4.5,-0.85) {bão hoà};
\node[font=\scriptsize, text=gray!60!black] at (axis cs:4.5,-0.85)  {bão hoà};

% ---- curves ----
\addplot[thick, aivnGreen]        {1/(1+exp(-x))};
\addlegendentry{$\sigma(z)$}
\addplot[thick, brandIndigo]      {tanh(x)};
\addlegendentry{$\tanh(z)$}
\addplot[thick, aivnGreen, dashed]   {(1/(1+exp(-x)))*(1-1/(1+exp(-x)))};
\addlegendentry{$\sigma'(z)$}
\addplot[thick, brandIndigo, dashed] {1-(tanh(x))^2};
\addlegendentry{$\tanh'(z)$}

% ---- the 0.25 ceiling ----
\draw[densely dotted, thick, red!70!black]
    (axis cs:-6,0.25) -- (axis cs:6,0.25);
\node[font=\scriptsize, text=red!70!black, anchor=west]
    at (axis cs:-5.9,0.38) {trần $\sigma' = 1/4$};
\end{axis}
\end{tikzpicture}

\begin{figurecaptionmodern}
\capfig{Sigmoid và $\tanh$ cùng đạo hàm của chúng trên $[-6,6]$. Vùng xám là miền
bão hoà, nơi đạo hàm tiến về $0$. Đường chấm đỏ là trần $\sigma' = 1/4$ mà đạo
hàm sigmoid không bao giờ vượt qua.\label{fig:sat}}
\end{figurecaptionmodern}
\end{figure}

\subsection{Nguyên nhân 2 --- bán kính phổ của ma trận trọng số}

Nguyên nhân thứ hai độc lập hoàn toàn với hàm kích hoạt: ngay cả với một mạng
tuyến tính ($\sigma = \mathrm{id}$, mọi $\sigma' = 1$), tích các ma trận trọng số
vẫn có thể suy biến.

Ký hiệu $\sigma_{\max}(\mathbf{A}) = \|\mathbf{A}\|_2$ là giá trị kỳ dị lớn nhất
(chuẩn phổ). Tính chất dưới nhân của chuẩn phổ,
$\|\mathbf{A}\mathbf{B}\|_2 \le \|\mathbf{A}\|_2\|\mathbf{B}\|_2$, cho ngay

\begin{infobox}{Chặn trên của tích ma trận}
\[
  \left\|\left(\prod_{k}\mathbf{W}^{[k]}\right)\mathbf{v}\right\|
  \;\le\;
  \left(\prod_{k}\sigma_{\max}\!\left(\mathbf{W}^{[k]}\right)\right)\|\mathbf{v}\| .
\]
\end{infobox}

Hệ quả là một phép phân loại sạch sẽ theo ba chế độ:

\begin{itemize}[leftmargin=*, itemsep=2pt]
  \item Nếu $\sigma_{\max}\!\left(\mathbf{W}^{[k]}\right) < 1$ với mọi $k$, tích
    các chặn là một cấp số nhân công bội nhỏ hơn $1$: gradient
    \textbf{suy giảm theo cấp số nhân}.
  \item Nếu $\sigma_{\max}\!\left(\mathbf{W}^{[k]}\right) > 1$ với mọi $k$, chặn
    trên phân kỳ. Trường hợp này không tự động kéo theo bùng nổ (chặn trên lớn
    không bắt buộc tích phải lớn), nhưng nếu thêm giả thiết các hướng kỳ dị
    trội không trực giao nhau --- điều xảy ra với xác suất cao khi trọng số khởi
    tạo ngẫu nhiên --- thì biên độ thực sự \textbf{tăng theo cấp số nhân}.
  \item Chế độ biên $\sigma_{\max} = 1$ là không bền: nhiễu số học bất kỳ cũng
    đẩy hệ sang một trong hai phía.
\end{itemize}

Đây chính là lý do \emph{khởi tạo} trọng số là một quyết định thiết kế chứ không
phải chi tiết kỹ thuật. Xavier~\cite{GlorotBengio2010} và He~\cite{He2015} không
làm gì khác ngoài việc chọn phương sai khởi tạo sao cho
$\sigma_{\max}(\mathbf{W})$ nằm gần $1$ ngay từ bước đầu tiên.

\subsection{Hợp nhất: Chặn trên bảo thủ của hệ số co giãn Jacobian}

Ghép hai nguyên nhân lại. Áp dụng tính dưới nhân của chuẩn phổ,
$\left\|\mathbf{A}\mathbf{B}\right\|_2 \le \left\|\mathbf{A}\right\|_2\left\|\mathbf{B}\right\|_2$,
cho từng thừa số của~\eqref{eq:central}:
\[
  \left\|\mathbf{M}^{[k]}\right\|_2
  \;\le\;
  \underbrace{\left\|\operatorname{diag}\!\left(\sigma'(\mathbf{z}^{[k-1]})\right)\right\|_2}_{\le\;\max\sigma'}
  \cdot
  \underbrace{\left\|\left(\mathbf{W}^{[k]}\right)^{\!\top}\right\|_2}_{=\;\sigma_{\max}(\mathbf{W}^{[k]})} .
\]
Do đó
\begin{equation}
  \left\|\boldsymbol{\delta}^{[l]}\right\|
  \;\le\;
  c^{\,L-l}\left\|\boldsymbol{\delta}^{[L]}\right\|,
  \qquad
  \boxed{\;c \;=\; \sigma_{\max}(\mathbf{W}) \cdot \max\sigma'\;}
  \label{eq:governing}
\end{equation}
với $c$ là \emph{chặn trên} của hệ số co giãn mỗi tầng, không phải giá trị thực
của nó.

\textbf{Vì sao chặn này bảo thủ.} Bất đẳng thức~\eqref{eq:governing} thay mỗi
thừa số bằng trường hợp xấu nhất của nó. Hệ số co giãn thực tế của một tầng,
$\left\|\mathbf{M}^{[k]}\mathbf{v}\right\| / \left\|\mathbf{v}\right\|$ với
$\mathbf{v}$ là vector sai số đang lan truyền ngược, còn phụ thuộc vào hai yếu tố
mà $c$ bỏ qua:
\begin{enumerate}[leftmargin=*, itemsep=1pt, label=(\alph*)]
  \item \textbf{Dữ liệu, qua tiền kích hoạt.} Ma trận
    $\operatorname{diag}(\sigma'(\mathbf{z}^{[k-1]}))$ được tính tại tiền kích
    hoạt thực $\mathbf{z}^{[k-1]}$, vốn phụ thuộc đầu vào $\mathbf{x}$ và mọi tầng
    phía trước. Ở vùng bão hoà của sigmoid, hoặc tại các nơ-ron ReLU chết
    ($\sigma'=0$), phần tử đường chéo nhỏ hơn hẳn $\max\sigma'$.
  \item \textbf{Sự thẳng hàng với các hướng kỳ dị.} Đẳng thức
    $\left\|(\mathbf{W}^{[k]})^{\top}\mathbf{v}\right\| = \sigma_{\max}(\mathbf{W}^{[k]})\left\|\mathbf{v}\right\|$
    chỉ xảy ra khi $\mathbf{v}$ nằm dọc vector kỳ dị trội. Một vector lan truyền
    ngược bất kỳ chỉ đặt một phần năng lượng lên các hướng đó, nên hệ số co giãn
    thực nằm đâu đó giữa $\sigma_{\min}$ và $\sigma_{\max}$.
\end{enumerate}

Vì vậy các kết luận rút ra được từ $c$ là kết luận \emph{một chiều}:
\begin{itemize}[leftmargin=*, itemsep=1pt]
  \item $c < 1$ là \emph{điều kiện đủ} để gradient co lại theo cấp số nhân: chặn
    trên $c^{\,L-l}\left\|\boldsymbol{\delta}^{[L]}\right\|$ đã giảm theo hàm mũ
    thì $\left\|\boldsymbol{\delta}^{[l]}\right\|$, bị chặn bởi nó, cũng phải giảm
    ít nhất nhanh như vậy.
  \item $c > 1$ chỉ làm cho bùng nổ \emph{có thể xảy ra}, không bảo đảm nó xảy
    ra: chặn trên tăng không kéo theo giá trị thực tăng, vì (a) và (b) có thể
    giữ hệ số co giãn thực dưới $1$.
  \item $c \approx 1$ chỉ bảo đảm chặn trên không co hay giãn theo hàm mũ; giá
    trị thực vẫn có thể suy giảm vì (a) và (b). Mục~\ref{sec:thucnghiem} đo
    trực tiếp khoảng cách giữa chặn và giá trị thực.
\end{itemize}

Bây giờ thay số cho sigmoid. Vì $\max\sigma' = 1/4$, để chặn trên đạt $c \ge 1$
ta cần
\[
  \sigma_{\max}(\mathbf{W}) \;\ge\; \frac{1}{\max\sigma'} \;=\; 4 .
\]
Nghĩa là ma trận trọng số phải có chuẩn phổ vượt $4$ \emph{chỉ để hoà vốn}. Nhưng
một ma trận với $\sigma_{\max} = 4$ khuếch đại tín hiệu \emph{tiến} lên bốn lần
mỗi tầng, đẩy các tiền kích hoạt $\mathbf{z}$ ra xa gốc --- đúng vào vùng bão
hoà, nơi $\sigma'$ sụp đổ về $0$ và giả thiết $\max\sigma' = 1/4$ không còn
đúng. Hệ thống tự phá huỷ điều kiện hoà vốn của chính nó.

\subsection{Nghịch lý phương sai: vì sao không thể phóng đại trọng số}

Lập luận ở cuối mục trước mới chỉ là một khẳng định định tính. Mục này chứng minh
nó: điều kiện $c\ge1$ \emph{không thể} đạt được bằng cách tăng thang đo trọng số,
vì chính việc tăng đó phá huỷ giả thiết mà nó dựa vào.

\subsubsection{Truy hồi phương sai theo chiều tiến}

Giả sử tại thời điểm khởi tạo:
\begin{enumerate}[leftmargin=*, itemsep=1pt, label=(G\arabic*)]
  \item $W^{[l]}_{ij}\sim\mathcal{N}(0, s^2)$ độc lập cùng phân phối, mọi tầng
    rộng $n$;
  \item các trọng số độc lập với các kích hoạt đầu vào của cùng tầng;
  \item các thành phần $a^{[l-1]}_j$ không tương quan với nhau;
  \item hàm kích hoạt hoạt động ở chế độ gần tuyến tính quanh gốc, tức
    $\sigma(z)\approx\sigma(0) + \sigma'(0)z$ với $\sigma'(0)$ là một hằng số bậc
    đơn vị.
\end{enumerate}

Đây là các giả thiết \emph{bậc một}, không phải đồng nhất thức. Chúng đúng tại
bước khởi tạo ngẫu nhiên và xấu đi khi huấn luyện tiến triển --- một điểm sẽ
quay lại ở Mục~\ref{sec:tongket} khi bàn về chuẩn hoá.

Với $z^{[l]}_i = \sum_{j=1}^{n} W^{[l]}_{ij}a^{[l-1]}_j$, do (G1)--(G3) và
$\mathbb{E}[W]=0$:
\[
  \operatorname{Var}\!\left(z^{[l]}_i\right)
  = \sum_{j=1}^{n}\operatorname{Var}\!\left(W^{[l]}_{ij}\right)
                  \operatorname{Var}\!\left(a^{[l-1]}_j\right)
  = n\,s^2\operatorname{Var}\!\left(a^{[l-1]}\right).
\]
Dùng (G4) để thay $\operatorname{Var}(a^{[l-1]})\approx\sigma'(0)^2
\operatorname{Var}(z^{[l-1]})$ và gộp hằng số $\sigma'(0)^2$ vào $s^2$, ta được
truy hồi
\begin{equation}
  \operatorname{Var}\!\left(\mathbf{z}^{[l]}\right)
  \;\approx\; \left(n s^2\right)\operatorname{Var}\!\left(\mathbf{z}^{[l-1]}\right)
  \;\Longrightarrow\;
  \operatorname{Var}\!\left(\mathbf{z}^{[l]}\right)
  \;\approx\; \left(n s^2\right)^{l}\operatorname{Var}\!\left(\mathbf{x}\right).
  \label{eq:varrec}
\end{equation}
Đại lượng $n s^2$ đóng vai trò hệ số khuếch đại phương sai mỗi tầng, hoàn toàn
song song với $c$ ở~\eqref{eq:governing} nhưng theo chiều ngược lại.

\subsubsection{Thế kẹt}

Với ma trận ngẫu nhiên $n\times n$ có phương sai phần tử $s^2$, định luật
Bai--Yin cho $\sigma_{\max}(\mathbf{W})\approx 2s\sqrt{n}$. Điều kiện hoà vốn
$\sigma_{\max}(\mathbf{W})\ge4$ của mục trước do đó tương đương
\[
  2s\sqrt{n}\;\ge\;4
  \quad\Longleftrightarrow\quad
  n s^2 \;\ge\; 4 .
\]
Thay vào~\eqref{eq:varrec}:
\[
  \operatorname{Var}\!\left(\mathbf{z}^{[l]}\right)\;\gtrsim\;4^{l}\operatorname{Var}(\mathbf{x}),
  \qquad\text{tức}\qquad
  \left|z^{[l]}\right| \;\sim\; 2^{\,l}\,\left|x\right| .
\]
Độ lớn tiền kích hoạt tăng theo cấp số nhân với chiều sâu. Nhưng ở vùng $|z|$
lớn, $\sigma(z)\to1$ và $1-\sigma(z)\approx e^{-|z|}$, nên
\begin{equation}
  \sigma'(z) = \sigma(z)\bigl(1-\sigma(z)\bigr) \;\approx\; e^{-|z|}
  \qquad (|z|\gg1).
  \label{eq:expdecay}
\end{equation}
Nghĩa là $\sigma'$ suy giảm \textbf{theo hàm mũ} của $|z|$, trong khi $|z|$ lại
tăng theo \emph{cấp số nhân} của $l$. Ghép hai điều đó:
\[
  \max\sigma'\big|_{\text{tầng }l}\;\sim\;e^{-2^{\,l}|x|},
\]
một sự sụp đổ nhanh hơn mọi thứ mà việc nâng $\sigma_{\max}(\mathbf{W})$ từ $1$
lên $4$ có thể bù lại --- vế sau chỉ là một thừa số hằng bằng $4$, vế trước là
hàm mũ kép theo chiều sâu.

\begin{notebox}
\textbf{Bùng nổ phương sai tiến ĐỔI LẤY triệt tiêu gradient ngược.} Với hàm kích
hoạt bão hoà, $c = \sigma_{\max}(\mathbf{W})\cdot\max\sigma'$ \emph{không thể}
được đưa về $1$ bằng cách chỉnh thang đo trọng số: mọi mức tăng $s$ đủ lớn để
nâng $\sigma_{\max}$ đều đẩy $\mathbf{z}$ ra vùng bão hoà theo~\eqref{eq:varrec},
nơi $\max\sigma'$ sụp đổ theo hàm mũ~\eqref{eq:expdecay}. Hai thừa số của $c$
không độc lập --- chúng bị ràng buộc ngược chiều qua chính thang đo trọng số.
\end{notebox}

Kết luận này có giá trị kiến tạo: nó \emph{loại trừ} toàn bộ một họ giải pháp và
chỉ ra chính xác ba lối thoát còn lại, mỗi lối phá vỡ một mắt xích của thế kẹt:

\begin{enumerate}[leftmargin=*, itemsep=2pt]
  \item \textbf{Đổi hàm kích hoạt} --- nâng $\max\sigma'$ lên $1$ mà không đụng
    tới $s$, phá vỡ~\eqref{eq:expdecay}. Đây là ReLU và các biến thể~\cite{He2015}.
  \item \textbf{Tái định tâm $\mathbf{z}$} --- giữ $\operatorname{Var}(\mathbf{z}^{[l]})$
    không đổi theo $l$ bất chấp~\eqref{eq:varrec}, tức cưỡng bức $ns^2\approx1$ ở
    mọi tầng và mọi bước huấn luyện. Đây là ý tưởng chung của khởi tạo
    Xavier~\cite{GlorotBengio2010} (đúng tại bước $0$) và của chuẩn
    hoá~\cite{IoffeSzegedy2015} (duy trì suốt quá trình).
  \item \textbf{Đổi cấu trúc Jacobian} --- thêm số hạng bậc không để tích thôi là
    tích thuần tuý, khiến câu hỏi về $c$ trở nên không còn quyết định. Đây là kết
    nối tắt~\cite{He2016ResNet}.
\end{enumerate}

Ba lối thoát này chính là ba cột đầu của bảng tổng kết ở Mục~\ref{sec:tongket};
chúng không phải một danh sách mẹo rời rạc mà là một phân loại vét cạn suy ra từ
\eqref{eq:governing} và~\eqref{eq:varrec}.

\begin{notebox}
\textbf{Triệt tiêu và bùng nổ không phải hai lỗi khác nhau.} Chúng là cùng một
cơ chế --- tích lặp của các Jacobian tầng --- quan sát ở hai phía của ngưỡng
$c = 1$. Mọi biện pháp khắc phục ở Mục~\ref{sec:tongket} đều là hệ quả của
\eqref{eq:governing}: hoặc nâng $\max\sigma'$, hoặc điều khiển
$\sigma_{\max}(\mathbf{W})$, hoặc giữ $\mathbf{z}$ ngoài vùng bão hoà, hoặc thay
đổi cấu trúc Jacobian để nó không còn là tích thuần tuý.
\end{notebox}
